\chapter{Detailed Design and Implementation}
\label{ch:design}
\pagestyle{fancy}

Following the theoretical framework and mathematical models presented in Chapter \ref{ch:analysis}, this chapter provides a comprehensive overview of the system's design and implementation. Detailed implementations regarding data processing, frontend, backend, and containerization are provided in order to illustrate the complete workflow, from the initial processing of patient data to the delivery of actionable insights through the user interface.

\section{Detailed Design}

As shown in Figure \ref{fig:system_design}, as a whole, the system design can be seen as an analytics pipeline which consists of data ingestion, storage in the form of structured relations and processing with multiple machine learning steps. The entire process seamlessly incorporates unsupervised methods, such as clustering and sequence mining, along with a Large Language Model in order to analyze raw data and present structured insights.

\begin{figure}[H]
\centering
\includegraphics[width=0.45\linewidth]{figs/chapter_5/system_design.png}
\caption{System Design Diagram}
\label{fig:system_design}
\end{figure}


\subsection{Synthea Data Ingestion and ETL Pipeline}
\label{ch:synthea_design}

The implementation of any of the algorithms described in this paper would not have been possible without the presence of an adequate dataset. An open-source Synthea model was chosen to serve that purpose due to its ability to provide realistic synthetic data on patients. Of the numerous files generated by the program, several administrative/financial datasets, namely \texttt{claims}, \texttt{claims\_transactions}, \texttt{payers}, \texttt{payer\_transitions}, \texttt{devices}, \texttt{imaging\_studies}, and \texttt{supplies}, were consciously omitted from the analysis process. The exclusion of these specific tables was done based on the fact that they address financial aspects of health care, which makes them irrelevant to our work.

\begin{figure}[H]
\centering
\includegraphics[width=0.8\linewidth]{figs/chapter_5/etl_pipeline.png}
\caption{ETL Pipeline Diagram}
\label{fig:etl_pipeline}
\end{figure}

\clearpage

Normalization is among the key concerns that go into designing the output relational database schema. Each table possesses its own well-defined primary key, and all non-key attributes are fully dependent on the entire key, not some part of it. Many-to-many relations are dealt with properly, employing junction tables, as exemplified by \texttt{patient} and \texttt{provider} connected via \texttt{patient\_provider} table.

Several calculated columns are provided in the database: age of patient is computed based on \texttt{birthdate}, duration of medication and careplan is derived from corresponding start and end date, \texttt{organization.utilization}, \texttt{provider.encounters}, and \texttt{provider.procedures} are counts. They remain in the database on purpose, in order to save the processing power on aggregating and joining data for queries, a common optimization technique for an analytical database in which reading is more frequent than writing.

\subsection{K-Prototypes Clustering Design}

Patient profiles and measurements include data of different types, both numeric (such as age) and categorical (such as gender and dominant care plan). In order to handle both data types natively without artificial encodings, the K-Prototypes algorithm was chosen. It extends K-Means by adding a categorical distance measurement to its cost function, as previously detailed in Section~\ref{ch:clustering_limitations}.

\subsubsection*{Selecting the Number of Clusters}

The optimal $k$ value was chosen by using the "Elbow" Plot method, based on the within cluster cost variation, for $k \in \{2, 3, \ldots, 9\}$. As seen in Figure~\ref{fig:ch5_elbow_plot}, a point of diminishing returns clearly appears when $k=4$. The reduction in cost when increasing $k$ from 3 to 4 was almost 900 points, while from 4 to 5 it was just about 350 points. This "elbow" point for $k=4$ for the chosen set of features is also supported by a benchmark that tunes various parameters for 63 different configurations, details of which are provided in Chapter~\ref{ch:testing}.

\begin{figure}[H]
\centering
\includegraphics[width=0.80\linewidth]{figs/chapter_5/elbow_plot.png}
\caption{Elbow plot used to select the number of clusters}
\label{fig:ch5_elbow_plot}
\end{figure}

\subsubsection*{Feature Engineering}

Many candidate features were created and tested for over 629 distinct runs of K-Prototypes algorithm before choosing the optimal features list. The set of features can be divided into two categories: \textbf{direct attributes} directly extracted from the database (e.g., age, gender), and \textbf{clinical signals} derived from the events in the raw table format. Feature derivation process is described by Figure~\ref{fig:ch5_feature_pipeline}. It takes raw data as an input, extracts features, passes high-cardinality and numeric variables to the bucketing layer, and finally forms a feature vector for K-Prototypes.

\begin{figure}[H]
\centering
\includegraphics[width=0.75\linewidth]{figs/chapter_5/feature_pipeline.png}
\caption{Feature engineering pipeline for the K-Prototypes input matrix}
\label{fig:ch5_feature_pipeline}
\end{figure}

The final configuration uses \textbf{5 numerical} and \textbf{8 categorical} features, which are described in detail below.

\paragraph{Numeric Features}

\begin{itemize}

    \item \textbf{age} -- is pulled straight from the patient medical record.

    \item \textbf{active\_conditions\_count} -- count of all those conditions that have not been resolved (no end date). Hence, it reflects the current chronic disease burden of a patient.

    \item \textbf{total\_medications} -- count of medication records associated with the patient at some point during their lifetime.

    \item \textbf{total\_procedures} -- count of procedures that have been conducted on the patient.

    \item \textbf{preventive\_to\_reactive\_ratio} -- number of wellness visits over the sum of acute visits and severe visits, with a unit added to the denominator to avoid division by zero for patients with no acute or severe encounters. The metric helps evaluate how proactive a patient is about their health: a value of 0 corresponds to a patient with no wellness visits, while higher values indicate a stronger focus on preventive care over reactive encounters.

\end{itemize}



\paragraph{Categorical Features}

\begin{itemize}

    \item \textbf{gender} -- pulled directly from the patient’s record.

    \item \textbf{has\_diabetes}, \textbf{has\_hypertension}, \textbf{has\_heart\_disease} -- binary attributes created through comparison with descriptions using clinical keywords. This triad of conditions was chosen based on their high prevalence and impact on care paths.

    \item \textbf{pregnancy\_history} -- a three-valued attribute (Active / Past / None) rather than a binary flag, since a pregnancy actively underway requires different medical attention than one that occurred previously.

    \item \textbf{condition\_burden} -- a binned version of \texttt{active\_conditions\_count}, discussed in more detail below.

    \item \textbf{bmi\_type} -- a binned form of average BMI calculated based on observations, detailed further below.

    \item \textbf{dominant\_careplan} -- the most severe care plan associated with the patient, binned from descriptions to clinical attributes, described in detail below.

\end{itemize}

\subsection{Pattern Mining Design}

The results from the Clustering process are in fact static representations of patients. They tell us which category the patients fit, but do not tell \emph{how} patients travel across the care system. In order to identify recurrent temporal patterns of clinical events, an additional layer of pattern mining was developed that would take advantage of the clustered patients. PrefixSpan algorithm (Section~\ref{ch:prefix_span_theory}), was selected for use in this case because of its efficient projection-based growth method in handling large vocabularies of hundreds of clinical event codes, making algorithms such as GSP infeasible.

The four-stage pipeline for the mining process can be seen in Figure~\ref{fig:ch5_pattern_pipeline}. At each stage, the implementation focuses on solving one unique challenge associated with the process. For instance, the process starts with organizing clinical data into patient sequences (Stage 1). Using the PrefixSpan approach, the data miner extracts frequent pattern candidates (Stage 2). In order to ensure the clinical significance of the findings, the candidates go through a process of ranking and scoring (Stage 3). Following this, a filtering process removes any duplicates (Stage 4), ensuring that the best results are obtained.

\begin{figure}[H]
\centering
\includegraphics[width=1\linewidth]{figs/chapter_5/pattern_pipeline.png}
\caption{Pattern mining pipeline stages}
\label{fig:ch5_pattern_pipeline}
\end{figure}

\subsubsection*{Sequence Construction}

The clinical events are generated from four different tables (\texttt{condition}, \texttt{medication}, \texttt{procedure}, \texttt{careplan}) and have been annotated with a type tag, thereby treating events with identical descriptions present in two different tables as two different clinical events. The process has been depicted in Fig.~\ref{fig:ch5_sequence_example} for a sample individual, where five raw rows from four different tables get condensed into one clinical sequence.

\begin{figure}[H]
\centering
\includegraphics[width=1\linewidth]{figs/chapter_5/sequence_example.png}
\caption{Example of a length 5 sequence}
\label{fig:ch5_sequence_example}
\end{figure}

\subsubsection*{Mining and Filtering}

\begin{figure}[H]
\centering
\includegraphics[width=0.7\linewidth]{figs/chapter_5/prefixspan_tree.png}
\caption{PrefixSpan projection tree}
\label{fig:ch5_prefixspan_tree}
\end{figure}

The PrefixSpan algorithm is applied on the generated clinical event sequences for the identification of the frequent paths. Figure~\ref{fig:ch5_prefixspan_tree} shows how the algorithm is carried out on a sample dataset of four sequences with minimum support of 2. Beginning from the empty prefix, the algorithm scans through the dataset and counts the occurrences of each item. The algorithm successfully identifies three frequent length 1 prefixes (\texttt{<C>}, \texttt{<M>}, and \texttt{<P>}, all with a support of 4), while immediately pruning infrequent items, such as \texttt{<CP>}, that fall below the required threshold.

After detecting these preliminary elements, PrefixSpan will recursively construct projected sub-databases for each of the frequent prefixes. For example, the projected sub-database $S|\texttt{<C>}$ contains only the suffixes of the sequences following the first occurrence of \texttt{C}. By processing this reduced-sized database, PrefixSpan can incrementally generate longer patterns and hence find \texttt{<C, M>} and \texttt{<C, P>} (support = 3). The process of projection followed by growth is performed recursively as can be seen from the generation of $S|\texttt{<C, M>}$ to yield a three-length pattern, \texttt{<C, M, P>}. This continues until there are no more frequent items to append. Since each projected sub-database is necessarily smaller than its predecessor, the search space will naturally become smaller and smaller.

One important feature of the architecture in this pipeline is the use of post-process clinical quality filters \emph{after} candidate generation and not through embedded in-loop constraints. With the use of a fully heuristic-independent approach in the pattern mining phase, the pipeline becomes very flexible, as clinical scoring and filtering can be optimized incrementally, without the overhead of re-executing the PrefixSpan algorithm.

\subsubsection*{Dual Scope Mining}

\begin{figure}[H]
\centering
\includegraphics[width=0.67\linewidth]{figs/chapter_5/sequence_modes.png}
\caption{Sequence Mining Operating Modes}
\label{fig:ch5_sequence_modes}
\end{figure}

In order to capture trends at both macro and micro levels (as shown in Figure~\ref{fig:ch5_sequence_modes}), the pattern mining process is performed in two modes. One mode is a global run that performs its process on the complete patient set, whereas another one is a cluster-specific run which runs the process on selected patients. The generated patterns are then indexed using their corresponding scopes, allowing for the incorporation of both global and cluster-level pathways in the same index structure. In addition to this, the outlier data is not ignored but is instead clustered separately.

\section{Detailed Implementation}

The Hospital Analytics Platform works as \textbf{Dockerized microservices}, ensuring that each calculation occurs in an isolated environment. The back end has \textbf{FastAPI} services in which the algorithmic layer provides \textbf{K-Prototypes Clustering} and \textbf{PrefixSpan}. An \textbf{AI Service} connects the \textbf{SQL Agent} and chat capabilities, transforming human language into queries and analysis, along with a \textbf{mail service} for transactional messages and \textbf{RabbitMQ} for communication between services. The front end is a \textbf{Next.js} application using REST APIs to present the results. The platform benefits its stakeholders: patients, doctors, and managers, by allowing them to analyze results from algorithms, manage patients' records, review dashboard statistics, and ask questions about the database using natural language through the \textbf{SQL Agent}. The rest of this paper will elaborate on the backend, frontend interface, \textbf{RabbitMQ} messaging, and Docker configuration.

\subsection{General Architecture}

\begin{figure}[H]
\centering
\includegraphics[width=0.90\linewidth]{figs/chapter_5/system_architecture.png}
\caption{System Architecture Diagram}
\label{fig:system_architecture}
\end{figure}

\subsubsection{Traefik}

\subsubsection{Backend Microservices}
\subsubsection*{Auth Service}
\subsubsection*{Users Service}
\subsubsection*{Analytics Service}
\subsubsection*{Clinical Core Service}
\subsubsection*{Queue Service}
\subsubsection*{Mail Service}

\subsubsection{Frontend Architecture}

\clearpage


\subsection{ETL Pipeline Implementation}

The Ingestion Script is an independent Python script (\texttt{db/etl\_script.py}), which consumes the CSV dataset generated by Synthea and dumps it into a normalized PostgreSQL database schema used by the other components of the system via SQLModel. The script operates independently of the FastAPI process, it runs only once for each database refresh operation and generates tables to be consumed by all subsequent layers. This script is part of the analytics microservice.

\subsubsection*{Connectivity}

Database connection is centralised in \texttt{db/db\_connector.py}, which provides an SQLAlchemy \texttt{engine}, a \texttt{SessionLocal} factory, a \texttt{get\_db()} dependency in FastAPI, and a \texttt{\_db\_retrieve(session, query)} function to retrieve data through raw queries as a list of dictionaries (\texttt{list[dict]}). Database credentials can be loaded using \texttt{python-dotenv} from the environment variables, making it possible to point the same codebase at different database instances across various deployment environments (such as local development, testing, and production) without requiring any source code modifications.

\subsubsection*{Source Dataset}

Synthea was run with a sufficiently large population (20{,}000 living patients) to generate a decades-long clinical history. The CSV files that come from Synthea to the ETL have the row counts seen below, of which only the eleven below will be ingested. All others including financial, imaging, and device records (\texttt{claims}, \texttt{claims\_transactions}, \texttt{payers}, \texttt{payer\_transitions}, \texttt{devices}, \texttt{imaging\_studies}, \texttt{supplies}) will be discarded on ingestion for the reasons explained in Section~\ref{ch:synthea_design}.

\begin{center}
\begin{tabular}{lr}
\hline
\textbf{Synthea CSV} & \textbf{Rows} \\
\hline
\texttt{patients.csv}        &        22,284 \\
\texttt{organizations.csv}   &         5,833 \\
\texttt{providers.csv}       &         5,833 \\
\texttt{encounters.csv}      &     1,211,526 \\
\texttt{allergies.csv}       &        21,761 \\
\texttt{procedures.csv}      &     3,225,114 \\
\texttt{medications.csv}     &       913,100 \\
\texttt{immunizations.csv}   &       335,814 \\
\texttt{careplans.csv}       &        70,197 \\
\texttt{conditions.csv}      &       727,501 \\
\texttt{observations.csv}    &    14,250,746 \\
\hline
\end{tabular}
\end{center}

\subsubsection*{Reducing Synthea to a Single Hospital}

The Synthea \texttt{organizations.csv} consists of 5,833 organisations in total covering an entire region (Massachusetts), with each organisation having its unique set of providers. The objective of this thesis is focused on a \emph{single} hospital spanning several years, rather than regional data. So it requires two filters to reduce the regional dataset to a single-hospital one.

\paragraph{One Organisation (Hospital).} The function \texttt{\_pick\_single\_org} reads the file \texttt{organizations.csv} and returns the row where the \texttt{UTILIZATION} column holds the highest value, since the specific row also represents the organization with the largest number of encounters in the bundle and therefore the richest longitudinal signal.

\begin{lstlisting}[style=pythonstyle]
def _pick_single_org() -> dict:
    df = pd.read_csv(dataset_path / "organizations.csv")
    row = df.loc[df["UTILIZATION"].idxmax()]
    return {"organization_id": row["Id"], "name": row["NAME"], ...}
\end{lstlisting}

\paragraph{Top 300 providers.} From the total pool of 5,833 providers, the top 300 are selected by ranking them on their encounter count in \texttt{encounters.csv}, since these 300 providers are the ones handling a clinically meaningful volume of activity.

\begin{lstlisting}[style=pythonstyle]
def _pick_top_provider_ids(max_providers: int = 300) -> list[str]:
    enc_df = pd.read_csv(dataset_path / "encounters.csv", usecols=["PROVIDER"])
    counts = enc_df["PROVIDER"].value_counts()
    return counts.head(max_providers).index.tolist()
\end{lstlisting}

\paragraph{Redistribution of encounters.} Blindly removing all encounters for the more than 5,500 omitted doctors would strip away most of the timeline. As such, all encounters for a physician whose data has been omitted are redistributed to one of the 300 retained providers. This occurs on a \emph{year-by-year} basis, cycling through the list of retained physicians sorted by their volume in the same year, ensuring that the temporal density of encounters per physician remains reasonable and that there is no sudden influx of patients seen decades prior to when the physician entered the dataset.

\begin{lstlisting}[style=pythonstyle]
# Build year -> ordered list of kept providers (most active first)
enc_year   = pd.to_datetime(df["START"], utc=True).dt.year
kept_mask  = df["PROVIDER"].isin(kept_ids)
year_provider_counts = (
    df[kept_mask].assign(year=enc_year[kept_mask])
                 .groupby(["year", "PROVIDER"]).size()
                 .reset_index(name="cnt")
                 .sort_values(["year", "cnt"], ascending=[True, False])
)
year_to_providers = {yr: grp["PROVIDER"].tolist()
                     for yr, grp in year_provider_counts.groupby("year")}
global_fallback   = year_provider_counts.groupby("PROVIDER")["cnt"].sum().idxmax()

# For each encounter from a dropped provider, pick the next kept provider
# from that year (round-robin via year_counters).
dropped_mask = ~encounters_df["provider_id"].isin(kept_ids)
years        = pd.to_datetime(encounters_df.loc[dropped_mask, "start_date"],
                              utc=True).dt.year
year_counters, new_providers = {}, []
for y in years:
    providers = year_to_providers.get(y, [global_fallback])
    idx       = year_counters.get(y, 0)
    year_counters[y] = idx + 1
    new_providers.append(providers[idx % len(providers)])
encounters_df.loc[dropped_mask, "provider_id"] = new_providers
\end{lstlisting}

\paragraph{Preservation of primary keys.} In the process of filtering and rewriting the data, the \emph{original} UUIDs for Synthea are preserved to serve as the primary keys (\texttt{patient\_id}, \texttt{provider\_id}, \texttt{organization\_id}, and \texttt{encounter\_id}). Since all child tables refer back to their corresponding parent tables using these primary key values, it was unnecessary to create a look-up table to translate the UUIDs.

\subsubsection*{CSV Ingestion}

The \texttt{pandas.read\_csv} function processes each Synthea CSV file by mapping it to only the columns required by the schema, converting all the column labels to lowercase, and performing minimal data cleaning before calling the \texttt{DataFrame.to\_sql} method:

\begin{lstlisting}[style=pythonstyle]
patients_df = df[list(keep_cols.keys())].rename(columns=keep_cols)
patients_df["first_name"] = patients_df["first_name"].apply(_strip_digits)
patients_df["last_name"]  = patients_df["last_name"].apply(_strip_digits)
patients_df["age"] = (
    patients_df["deathdate"].fillna(pd.Timestamp.now()) - patients_df["birthdate"]
).dt.days // 365

patients_df.to_sql(name="patient", con=engine, if_exists="append", index=False)
\end{lstlisting}

\noindent The synthetic IDs that Synthea appends to first and last names (\textit{e.g.,}\ \texttt{Maria123}) are stripped out using a regular expression helper function since we want the data display to appear natural. \texttt{to\_sql} is called with \texttt{if\_exists="append"} and (for the large tables) with \texttt{method="multi"}, allowing \texttt{psycopg2-binary} to bundle the rows into a single multi-value \texttt{INSERT} statement; this dramatically decreases the loading time of the entire Synthea dataset (approximately 20{,}000 patients and several million rows across the child tables) from tens of minutes to a few minutes on a developer machine.

\subsubsection*{Many-to-many Relationships}

Tables that represent a many-to-many relationship - the patient/doctor relationship being an obvious example - are filled by running a second scan of the \texttt{encounter} table that produces the junction table using \texttt{DISTINCT} on the tuple:

\begin{lstlisting}[style=sqlstyle]
INSERT INTO patient_provider (patient_id, provider_id)
SELECT DISTINCT patient_id, provider_id FROM encounter
WHERE provider_id IS NOT NULL;
\end{lstlisting}

\subsubsection*{Execution Order}

The table-specific load processes are wired together in \texttt{run\_etl}, which enforces the dependency of parents before children for the foreign keys, namely the patients, one organisation, and 300 providers must exist prior to any encounter records being loaded, which must be loaded before any allergies, procedures, medications, immunisations, careplans, conditions, or observations.

\begin{lstlisting}[style=pythonstyle]
def run_etl(engine):
    _truncate_all_tables(engine)
    load_patients(engine)
    load_organizations(engine)
    load_providers(engine)
    load_encounters(engine)
    load_patient_providers(engine)
    load_allergies(engine)
    load_procedures(engine)
    load_medications(engine)
    load_immunizations(engine)
    load_careplans(engine)
    load_conditions(engine)
    load_observations(engine)
\end{lstlisting}

In contrast, \texttt{\_truncate\_all\_tables} is executed in the \emph{reverse} order to ensure that all \texttt{TRUNCATE \dots CASCADE} operations will not violate any foreign keys.

\subsubsection*{Computed Columns}

The materialisation of computed columns mentioned in Section~\ref{ch:synthea_design} happens via plain SQL after data loading, which consists of issuing \texttt{ALTER TABLE \dots ADD COLUMN}, followed by \texttt{UPDATE \dots FROM}. For instance, the age of the patient is recalculated based on the real-time system clock so the attribute will always be up-to-date:

\begin{lstlisting}[style=sqlstyle]
UPDATE patient
SET age = DATE_PART('year', AGE(COALESCE(deathdate, CURRENT_DATE), birthdate));
\end{lstlisting}

and counts like \texttt{provider.encounters} have their data populated with a correlated aggregation update to ensure that all subsequent analysis queries skip the join operation. The foreign key constraints are defined at the end of the script using \texttt{ALTER TABLE \dots ADD CONSTRAINT}, ensuring that the constraints are defined after loading for performance reasons.

\subsubsection*{Libraries}

The principal libraries involved in the ETL layer are \texttt{pandas}, which is responsible for CSV processing, \texttt{SQLAlchemy} and \texttt{SQLModel} for schema definition and connection pooling, respectively, \texttt{psycopg2-binary}, which serves as the PostgreSQL driver, and finally \texttt{python-dotenv}, which enables environment variable loading. Bulk loads do not use ORM insertions; all the heavy writes happen through \texttt{DataFrame.to\_sql} or raw \texttt{text()} statements.

\subsection{K-Prototypes Pipeline Implementation}

The clustering pipeline consists of a script file called \texttt{k\_prototypes.py}. This script can be executed using \texttt{python k\_prototypes.py} in the appropriate directory and is not connected to FastAPI since it is time-consuming and should be done offline. The function for clustering is located in the \texttt{kmodes} at \texttt{kmodes.kprototypes.KPrototypes}. Everything else (feature extraction, scaling, projection, evaluating, plotting, persisting, and summarising with AI) is part of the wrapper around this algorithm. The entire wrapper is part of the \texttt{ml-batch} microservice.

\subsubsection*{Feature Construction}
\label{ch:feature_construction}

The function \texttt{build\_patient\_features(engine)} invokes an array of standalone SQL queries using \texttt{pandas.read\_sql(text(\dots), engine)}. Each query retrieves a compact dataframe, indexed by \texttt{patient\_id}, for each distinct feature category (patient demographics, encounters, condition flags, average BMI/BP/glucose, care plan, physicians, off-hours flag, substance abuse, pain level, readmission level, and so forth). The function performs a chain of \texttt{df.merge(right\_df, on="patient\_id", how="left")} operations on those dataframes, and then applies type-aware imputation.

Here is an example of the SQL magic, illustrating the process used to derive the \emph{dominant care plan} categorical variable (which constitutes one of the cluster-determining features). The 40 care plan categories are bucketized into 7 clinical categories based on severity rank with \texttt{ROW\_NUMBER() OVER (PARTITION BY patient\_id ORDER BY severity\_rank ASC)}.

\begin{lstlisting}[style=sqlstyle]
WITH CategorizedPlans AS (
    SELECT patient_id,
           CASE
               WHEN description IN (
                   'Cancer care plan (record artifact)',
                   'Dialysis care plan (record artifact)',
                   'Terminal care (regime/therapy)'
               ) THEN 'Complex, Oncology & Palliative'
               WHEN description = 'Routine antenatal care (regime/therapy)'
                                 THEN 'Maternal & Antenatal'
               WHEN LOWER(description) ~ '(fracture|wound|surgery|injury|trauma)'
                                 THEN 'Acute Trauma, Wound & Surgery'
               WHEN LOWER(description) ~ '(respiratory|asthma|pulmonary|covid)'
                                 THEN 'Respiratory & Infectious'
               WHEN LOWER(description) ~ '(diabetes|heart|hypertension|cardiac)'
                                 THEN 'Chronic Cardio-Metabolic'
               ELSE 'General Rehab & Routine'
           END AS care_plan_category,
           CASE
               WHEN description IN (
                   'Cancer care plan (record artifact)',
                   'Dialysis care plan (record artifact)',
                   'Terminal care (regime/therapy)'
               ) THEN 1
               WHEN description = 'Routine antenatal care (regime/therapy)' THEN 2
               WHEN LOWER(description) ~ '(fracture|wound|surgery|injury)'   THEN 3
               WHEN LOWER(description) ~ '(respiratory|asthma|pulmonary)'    THEN 4
               WHEN LOWER(description) ~ '(diabetes|heart|hypertension)'     THEN 5
               ELSE 6
           END AS severity_rank
    FROM careplan
),
RankedPatientPlans AS (
    SELECT patient_id, care_plan_category,
           ROW_NUMBER() OVER (
               PARTITION BY patient_id ORDER BY severity_rank ASC
           ) AS rank_id
    FROM CategorizedPlans
)
SELECT patient_id, care_plan_category AS dominant_careplan
FROM RankedPatientPlans
WHERE rank_id = 1;
\end{lstlisting}

The same technique is used for the other derived categories as well: \texttt{substance\_use\_profile}, \texttt{pain\_management\_tier}, \texttt{readmission\_risk\_tier}, \texttt{cholesterol\_risk\_category}, \texttt{respiratory\_disease\_profile}, and many more.

\subsubsection*{Feature Bucketing}

After the feature frames are merged, the preprocessing steps utilize an intentional approach for bucketing and imputation for dealing with the sparsity of clinical features:

\begin{lstlisting}[style=pythonstyle]
for col in ["wellness_count", "inpatient_count", ...]:
    df[col] = df[col].fillna(0)        # zero is semantically correct
df["dominant_careplan"] = df["dominant_careplan"].fillna("None")
df["bmi_type"] = pd.cut(df["avg_bmi"],
                        bins=[-inf, 18.5, 25, 30, 35, 40, inf],
                        labels=["Underweight", "Normal", "Overweight",
                                "Obese I", "Obese II", "Obese III"]) \
                  .astype(object).where(df["avg_bmi"].notna(), "Unknown")
\end{lstlisting}

\noindent Numerical features that are sparse, such as body mass index, are intentionally bucketed with \texttt{pd.cut} rather than being filled with their mean value. Given that K-Prototypes assumes that numerical centroids represent real-world measurements, adding the mean artificially would result in clusters being affected. These are then grouped into clinical buckets and any missing information is sent to a \texttt{"Unknown"} group. On the other hand, count-type features can be simply imputed with zeros since this accurately represents a lack of events.

\subsubsection*{The Role of \texttt{log1p}}

\texttt{total\_\allowbreak medications}, \texttt{lifetime\_\allowbreak medication\_\allowbreak days}, and \texttt{total\_\allowbreak encounter\_\allowbreak duration\_\allowbreak hours} have strong right tails: a few cases with many medications or long durations have values an order of magnitude above the median. For those columns, there are two steps in the pipeline where \texttt{np.log1p} is used in order to reduce such extreme values:

\begin{lstlisting}[style=pythonstyle]
# inside remove_outliers_zscore; z-scores on raw skewed counts would
# flag a quarter of the cohort as outliers purely on distribution shape.
for i, col in enumerate(cols):
    if col in _SKEWED_OUTLIER_COLS:
        data[:, i] = np.log1p(np.clip(data[:, i], 0, None))

# inside run_kprototypes, so the scaler does not push the "normal" mass
# of patients into a tiny corner of [0, 1].
for col in ["total_medications", "lifetime_medication_days",
            "total_encounter_duration_hours"]:
    if col in numerical_cols:
        df_scaled[col] = np.log1p(df_scaled[col].clip(lower=0).astype(float))
\end{lstlisting}

\noindent \texttt{log1p} (\(\log(1+x)\)) is better than simply taking \(\log(x)\) because \texttt{log1p} is defined for \(x=0\) (and the mode of these features is zero). This transformation is done only on \texttt{df\_scaled} (the input matrix for the K-Prototypes algorithm) and \emph{not} on \texttt{df\_work} (the frame returned for the LLM prompt and the database).

\subsubsection*{PCA for Severity Compression, FAMD for Visualisation}

There are two distinct dimensionality reduction methods applied for two very different applications. In broad terms, the dimensionality reduction methods, PCA, and FAMD perform two distinct roles: Feature Compression and Visualization. On a fundamental level, the methods of feature compression transform an original set of features that describe a patient through a mathematical process to represent him/her in a low-dimensional space. With careful consideration to maintain maximum variability and correlation between the variables, this technique reduces an extremely complicated set of features describing a patient to a single $(x,y)$ coordinate pair on a 2D plane.

\vspace{0.5em}
\textbf{PCA} (principal component analysis) accomplishes this task for continuous variables by determining an orthogonal basis for the input space that orders axes based on their variance contribution. In this case, applied to the three correlated clinical severity scores, the first principal component contains the common "how sick is this patient" information in one variable that can be weighted against other numericals by the \textbf{K-Prototypes} algorithm. As a safeguard, configurations containing \texttt{clinical\_severity\_score} and any of its components are disallowed, as this would be redundant variance.

\paragraph{Principal component analysis for the three components representing severity.} If the \texttt{clinical\_severity\_score} composite feature is present in the configuration, it is obtained by applying the \texttt{log1p} function to the three raw severity features (\texttt{total\_procedures}, \texttt{total\_medications}, \texttt{active\_conditions\_count}), followed by standardisation, and then projecting them onto the first principal component. Since the resulting axis is already a z-score by construction, the outer min-max scaler skips it:

\begin{lstlisting}[style=pythonstyle]
sev_input        = np.log1p(df[sev_cols].clip(lower=0).values.astype(float))
sev_standardized = StandardScaler().fit_transform(sev_input)
sev_pca          = PCA(n_components=1)
css              = sev_pca.fit_transform(sev_standardized)[:, 0]
print(f"explained variance ratio: {sev_pca.explained_variance_ratio_[0]:.4f}")

cols_to_scale = [c for c in numerical_cols if c != "clinical_severity_score"]
df_scaled[cols_to_scale] = scaler.fit_transform(df_scaled[cols_to_scale].values)
\end{lstlisting}

\paragraph{FAMD for the 2-D visualisation.} The scatter plot used in the dashboard is not computed using PCA but rather \textbf{FAMD} (Factor Analysis of Mixed Data), which comes from the \texttt{prince} library. This approach is simply the natural extension of PCA to the case where there is a mixture of quantitative and qualitative data types; quantitative data is standardised, while qualitative data is one-hot-encoded and then reweighted based on their inverse category frequencies.

\begin{lstlisting}[style=pythonstyle]
import prince
famd   = prince.FAMD(n_components=2, random_state=42).fit(df_feat)
coords = famd.transform(df_feat).values  # shape (n, 2)
\end{lstlisting}

PCA can be applied to numerics only; applying an ordinal encoding to categoricals prior to PCA destroys the categorical distance geometry utilized by K-Prototypes in its clustering process. By giving both types of variables equal treatment, FAMD creates a projection that is faithful to the space being clustered into.

\subsubsection*{Scaling the Numerical Block}

Three numerical scalers are exposed through a name-to-instance map and chosen per run. Each one is taken straight from \texttt{sklearn.preprocessing} and rewrites the column with a different rule:

\begin{lstlisting}[style=pythonstyle]
_SCALER_MAP = {
    "minmax":   MinMaxScaler(),    # x' = (x - min) / (max - min) -> [0, 1]
    "robust":   RobustScaler(),    # x' = (x - median) / IQR      -> centered on 0, unbounded
    "standard": StandardScaler(),  # x' = (x - mean)   / std      -> ~ N(0, 1), unbounded
}
\end{lstlisting}

\paragraph{MinMaxScaler.} Squeezes each column linearly into the finite interval \([0, 1]\). The minimum value maps to zero, the maximum to one, and all other values map linearly within that range. The shape of the distribution is preserved exactly, with only the support changing.

\paragraph{StandardScaler.} Centers each column around its own mean and divides by its own standard deviation, resulting in a column whose approximate distribution is \(\mathcal{N}(0, 1)\). The centering and unit variance are quite useful in practice, but the unbounded nature of the Gaussian causes the heavy tails in the column to remain unaltered.

\paragraph{RobustScaler.} Centers around the median and divides by the inter-quartile range (IQR). Both of these measures are robust to outliers, making this the best scaler to use when there are any real outliers present in the data; however it is still an unbounded output.

\paragraph{Why MinMax works best for K-Prototypes.} The K-Prototypes objective consists of a squared-Euclidean distance between the numerical features and their corresponding cluster prototypes, plus a \(\gamma\)-scaled Hamming distance between the categorical features and their corresponding cluster prototypes, where the Hamming term is bounded in \([0, n_{\text{cat}}]\). By scaling the numerical features to the same \([0, 1]\) interval using \texttt{MinMaxScaler}, the two terms have similar magnitudes no matter how far apart the numbers are from each other, hence \(\gamma\) can be automatically tuned by \texttt{kmodes} without any one term overpowering the other. Although \texttt{StandardScaler} centers at zero, it does not eliminate the long tails in the distribution, which increases the importance of outlier-heavy features for the numerical objective. \texttt{RobustScaler}, while more robust to outliers, still has a few outlying samples with extremely high or low numerical values (not within \([0, 1]\)), which resulted in less stable clusters on average across the benchmark.

\subsubsection*{Z-score Outlier Detection}
\label{ch: z_score}

Before the process of clustering begins, function \texttt{remove\_outliers\_zscore()} scans each patient according to a static list of numerical attributes. Z-Score of any value \(x\) means how far away \(x\) is from the average of the respective column measured by standard deviations. Thus, a patient with \(z=6\) on the \texttt{lifetime\_medication\_days} feature will be six standard deviations greater than the average value in this column. Outliers' scan goes over the original numerical \texttt{numpy} array of the patient data frame (NOT over the output of cluster analysis - since there are no clusters formed yet), but with an initial \texttt{log1p()} transformation applied to those columns with strong right skewness.

\begin{lstlisting}[style=pythonstyle]
data = df[cols].to_numpy(dtype=float)
for i, col in enumerate(cols):
    if col in _SKEWED_OUTLIER_COLS:
        data[:, i] = np.log1p(np.clip(data[:, i], 0, None))

mask = np.ones(len(df), dtype=bool)
for i, col in enumerate(cols):
    z     = np.abs((data[:, i] - np.nanmean(data[:, i])) / np.nanstd(data[:, i]))
    mask &= z < threshold      # threshold = 6 in production
\end{lstlisting}

The patient is retained iff \emph{all} columns' z-scores remain under the threshold. Notably, outliers are not removed, because the function produces two frames, and the \texttt{df\_outliers} one is propagated throughout the processing chain and receives the label value \(-1\) upon persistence. The result is that the outliers form a distinct cluster on the dashboard (\emph{outliers}), are explained by the LLM using the same prompt as a normal cluster, and even mined separately as a unique scope by PrefixSpan. No patient is thereby ignored in analysis; at the same time, the centroid for each of the four actual clusters represents an ordinary patient, not outliers.

\subsubsection*{Running the Algorithm}

The clustering call will build a \texttt{numpy} matrix where numeric columns appear before string columns; the latter will be stored as strings and their column indices will be passed through \texttt{categorical=}. It’s important to note that constructing the matrix piece by piece ensures that the Cython loop inside \texttt{kmodes} runs quickly:

\begin{lstlisting}[style=pythonstyle]
num_matrix    = df_scaled[numerical_cols].to_numpy(dtype=np.float64)
cat_matrix    = df_scaled[categorical_cols].astype(str).to_numpy()
matrix        = np.concatenate([num_matrix, cat_matrix], axis=1)
cat_positions = list(range(len(numerical_cols),
                           len(numerical_cols) + len(categorical_cols)))
kproto = KPrototypes(n_clusters=k, init="Huang",
                     random_state=rs, n_jobs=-1, gamma=gamma)
clusters = kproto.fit_predict(matrix, categorical=cat_positions)
\end{lstlisting}

\texttt{init="Huang"} is the frequency-based initialization scheme from Huang (1998), which seeds prototypes using attribute frequency distributions, and combined with a fixed \texttt{random\_state} ensures reproducibility across runs. If \texttt{gamma=None}, \texttt{kmodes} computes \(\gamma\) by using the standard deviation of the (previously min-max normalized) numeric portion, ensuring that both parts of the objective function are balanced.

\subsubsection*{Configuration Loop and Persistence to the Database}

A list of feature configurations is stored in \texttt{configs.py} using the format \texttt{CONFIGS = [\{"version": \dots, \allowbreak "preferred\_\allowbreak k": k, \allowbreak "numerical": [\dots], \allowbreak "categorical": [\dots]\}]}. \textbf{63 configurations} (\texttt{v1} to \texttt{v63}) have been created during the experimentation process, where each configuration tests a new hypothesis (such as severity-only, utilization-only, behaviour-focused, cardio-metabolic, demographic baseline only, and many more - the file archives all retired configurations with comments for an audit trail of the search process). The complete benchmarking for all 63 configurations and the reasoning behind choosing \texttt{v55} as the champion model is explained in Chapter~\ref{ch:testing}; the selected configuration remains uncommented in the code:

\begin{lstlisting}[style=pythonstyle]
{
    "version": 55,
    "preferred_k": 4,
    "numerical": [
        "age",
        "active_conditions_count",
        "total_medications",
        "total_procedures",
        "preventive_to_reactive_ratio",
    ],
    "categorical": [
        "gender",
        "pregnancy_history",
        "has_heart_disease",
        "has_diabetes",
        "has_hypertension",
        "condition_burden",
        "bmi_type",
        "dominant_careplan",
    ],
},
\end{lstlisting}

The primary loop of \texttt{k\_prototypes.py} runs over the Cartesian product of \texttt{THRESHOLDS}, \texttt{CONFIGS}, \texttt{SCALERS}, \texttt{GAMMAS}, \texttt{RANDOM\_STATES}, and the selected \(k\)s. Even when the production configuration forces all axes to have only one value each, this is the construct that makes parameter tuning possible, simply replace any of those singletons with a non-singleton list and you will produce a benchmark sweep:

\begin{lstlisting}[style=pythonstyle]
for threshold in THRESHOLDS:                          # outlier sensitivity
    df_cleaned, df_outliers = remove_outliers_zscore(df, threshold=threshold)
    for cfg in CONFIGS:                               # 63 versions
        for scaler_name in SCALERS:                   # minmax / robust / standard
            for gamma in GAMMAS:                      # auto or fixed
                for rs in RANDOM_STATES:              # init seeds
                    for k in k_values:                # cluster count
                        df_result, X_num_scaled, clusters, _, _ = run_kprototypes(...)
                        report, eval_metrics = full_evaluation_report(...)
                        # one row per run, written incrementally
                        general_rows.append({"version": ..., **eval_metrics})
                        flush_general_conclusion()
\end{lstlisting}

\noindent Each iteration produces its own \texttt{conclusions.txt}, \texttt{conclusions.xlsx}, FAMD scatter plot, heatmap, bar plot, and stacked categorical distribution inside a folder named for the particular iteration inside \texttt{outputs/}; however, a unique \texttt{general\_conclusion.xlsx} combines all the metrics from each run in order to make comparison easier. Most importantly, all metrics used for evaluation (silhouette, Gower silhouette, Calinski--Harabasz, and Davies--Bouldin, both in raw numerical space and in the FAMD projection) are recalculated each time with \texttt{full\_evaluation\_report}, which ensures that there is a complete signature for each run in the benchmark Excel file. The detailed result analysis and selection of \texttt{v55} are presented in Chapter~\ref{ch:testing}.

\subsubsection*{From DataFrame to \texttt{numpy} and Back to Words}

The K-Prototypes algorithm and the sklearn metrics use \texttt{numpy} arrays, but the prompts from the LLM model, the cards from the dashboard, and the \texttt{cluster\_\allowbreak patient.features} JSONB field require the original strings to be restored (\texttt{gender = "F"}, \texttt{bmi\_\allowbreak type = "Obese II"}, \texttt{dominant\_\allowbreak careplan = "Chronic Cardio-Metabolic"}). This is achieved by having \texttt{run\_\allowbreak kprototypes} maintain two frames simultaneously, one \emph{raw} (\texttt{df\_\allowbreak work}), which is returned to the persistence layer and LLM prompt, and the other (\texttt{df\_\allowbreak scaled}) containing only the transformations required for clustering (log1p, severity PCA, scaler), which is converted to the \texttt{numpy} array. The integer array of cluster labels from the \texttt{fit\_\allowbreak predict} call is then appended to the raw frame as a single column; no further "decoding" of the results is necessary.

\begin{lstlisting}[style=pythonstyle]
df_work   = df.copy()        # raw values, go to the LLM and the DB
df_scaled = df_work.copy()   # transformed copy, goes to numpy / sklearn

# ... log1p, severity PCA, scaler.fit_transform applied to df_scaled only ...

clusters = kproto.fit_predict(matrix, categorical=cat_positions)  # ndarray[int]
df_work["cluster"] = clusters    # one extra column, original strings intact
\end{lstlisting}


\subsubsection*{Reuse of Derived Features Across Configurations}

The 21 SQL features generated by \texttt{build\_patient\_features} are calculated \emph{only once} for each run and stored within the \texttt{df} dataframe. In the configuration loop, we simply \emph{select} the columns we need from that dataframe; no SQL will be executed when going from one configuration \(v_i\) to the next \(v_{i+1}\). That's why 63-config sweeping is possible - the computation of heavy lifting (extracting the rows, LEFT JOIN, bucketing categoricals) is done once, and then K-Prototypes fitting becomes an in-memory numpy calculation on the very same frame.

The same logic is employed on the persistence end: the \texttt{save\_cluster\_results} function is passed \texttt{extra\_feature\_cols=[...]}, so all the columns that aren't used in the clustering (like \texttt{avg\_bmi}, \texttt{qaly}, \texttt{has\_cancer}, and \texttt{smoking\_status}) are saved within the \texttt{cluster\_patient.features} JSONB column. The cluster detail and patient detail views in the frontend read the JSONB, which means the non-used-by-clustering features will be saved for use with the analytical API and summaries for each AI prompt.

\subsubsection*{AI-driven Cluster Cards on the Frontend}

Once the clustering step is done, \texttt{get\_clusters\_summary} generates a prompt that combines information about each cluster’s numeric mean, categorical composition, top conditions, medications, procedures, care plan distribution, and population context into one string and sends it via \texttt{ask\_ai} (same OpenAI-powered function the SQL agent uses). The prompt requests the following from the model:

\begin{lstlisting}[style=pythonstyle]
[Cluster ID] | [#HexColor] | [Short Title] | [1-2 sentence Description]
            | [1 short Insight] | [1 short Recommendation]

Choose a suggestive hex color for each cluster that reflects its clinical
profile (e.g. red tones for high-risk, green for healthy, blue for elderly
chronic, orange for moderate, purple for psychiatric, teal for maternal).
The outlier cluster (-1) should always be gray (#9CA3AF).
\end{lstlisting}

In order for the model to have sufficient context with which to select a color, a title, and a pertinent clinical insight, all features calculated for each cluster in the previous step are concatenated to the prompt in their verbatim form. For each cluster, six lines are added: the important average values (\texttt{age}, \texttt{qaly}, \texttt{daly}, \texttt{total\_procedures}, \texttt{total\_medications}, \texttt{active\_\allowbreak conditions\_\allowbreak count}, \texttt{lifetime\_\allowbreak medication\_\allowbreak days}, \texttt{total\_\allowbreak encounter\_\allowbreak duration\_\allowbreak hours}), the demographics and disease-prevalence dictionary (mode gender, mode marital status, mode encounter class, mode care plan, prevalence of diabetes, hypertension, heart disease, cancer, polypharmacy, BMI distribution, BP distribution, glucose distribution, substance use distribution, readmission distribution), the careplan share, and the top conditions, medications, and procedures pulled from the raw clinical tables joined onto the cluster labels:

\begin{lstlisting}[style=pythonstyle]
prompt = f"""We clustered {len(df_clustered)} patients into groups, plus
{len(df_outliers)} outlier patients. The outliers are labeled as Cluster -1.

Your task is to write extremely brief text to fit inside small frontend UI cards.
... (formatting + colour palette rules omitted) ...

=========================================
CLUSTER DATA SUMMARIES:
=========================================
"""
for c in clusters:
    c_df   = df_all[df_all['cluster'] == c]
    c_size = len(c_df)
    c_pct  = (c_size / total_patients) * 100

    c_means = c_df[['age', 'qaly', 'daly',
                    'total_procedures', 'total_medications',
                    'active_conditions_count', 'lifetime_medication_days',
                    'total_encounter_duration_hours']].mean().round(1).to_dict()

    c_cats = {
        'dominant_gender':    c_df['gender'].mode()[0],
        'dominant_marital':   c_df['marital'].mode()[0],
        'main_encounter_type': c_df['dominant_encounter_class'].mode()[0],
        'dominant_careplan':  c_df['dominant_careplan'].mode()[0],
        'diabetes_rate':      f"{(c_df['has_diabetes']      == 'Yes').mean()*100:.0f}%",
        'hypertension_rate':  f"{(c_df['has_hypertension']  == 'Yes').mean()*100:.0f}%",
        'heart_disease_rate': f"{(c_df['has_heart_disease'] == 'Yes').mean()*100:.0f}%",
        'cancer_rate':        f"{(c_df['has_cancer']        == 'Yes').mean()*100:.0f}%",
        'polypharmacy_rate':  f"{(c_df['is_polypharmacy']   == 'Yes').mean()*100:.0f}%",
        'bmi_mix':            c_df['bmi_type'           ].value_counts(normalize=True).round(2).to_dict(),
        'bp_mix':             c_df['bp_risk_category'   ].value_counts(normalize=True).round(2).to_dict(),
        'glucose_mix':        c_df['glucose_risk_category'].value_counts(normalize=True).round(2).to_dict(),
        'substance_use_mix':  c_df['substance_use_profile'].value_counts(normalize=True).round(2).to_dict(),
        'readmission_mix':    c_df['readmission_risk_tier'].value_counts(normalize=True).round(2).to_dict(),
    }
    careplan_mix = c_df['dominant_careplan'].value_counts(normalize=True).round(2).head(4).to_dict()
    c_cond = top_conditions [top_conditions ['cluster'] == c]['description'].head(4).tolist()
    c_meds = top_medications[top_medications['cluster'] == c]['description'].head(3).tolist()
    c_proc = top_procedures [top_procedures ['cluster'] == c]['description'].head(3).tolist()

    prompt += f"\n### DATA FOR CLUSTER {c} ({c_size} patients, {c_pct:.1f}% of total) ###\n"
    prompt += f"- Key Averages: {c_means}\n"
    prompt += f"- Demographics & Disease Rates: {c_cats}\n"
    prompt += f"- Careplan Mix (share of patients): {careplan_mix}\n"
    prompt += f"- Top Conditions:  {c_cond}\n"
    prompt += f"- Top Medications: {c_meds}\n"
    prompt += f"- Top Procedures:  {c_proc}\n"
\end{lstlisting}

The reason to include the entire dictionary, with all values of \texttt{c\_means}, \texttt{c\_cats}, and \texttt{careplan\_mix} dictionaries in the prompt, is the fact that the LLM must not speculate on anything but the information it gets. The exact numerical measures a clinician would give when describing the cluster, in the same exact units (\texttt{"Obese II"}, \texttt{"Chronic Cardio-Metabolic"}, \texttt{73\%}), are used as input so that the system generates a fitting color and a specific cluster name (red for a cardiometabolic cluster, teal for a maternal one, etc.). The output from the model is then parsed by \texttt{\_parse\_cluster\_summaries} (a simple regex for cluster ID within brackets followed by \texttt{split('|')}), and the result is saved literally into \texttt{cluster} table - that's the exact data you see in cluster summaries on the dashboard.

The same per-cluster values (average age, average BMI, systolic BP, glucose level, QALY/DALY, number of active diagnoses, medications, procedures, days since last encounter, gender balance, smoker distribution, dominant care plan and encounter class, and percentages of diabetes, hypertension, heart disease, cancer, and polypharmacy patients) are requested by the other prompt, the on-demand prompt \texttt{PROMPT\_AI\_ANALYSIS} defined in \texttt{cluster\_service.py}, that will be triggered on the first visit of the detailed cluster page by a user. Compared to the dashboard-card prompt, the novelty here is the inclusion of each metric alongside its population average, allowing the LLM to describe each deviation in relative terms (e.g., "higher/lower than the average"), and the injection of the silhouette and Davies--Bouldin scores available via \texttt{cluster\_meta.metrics} so the model has an idea about the credibility of the cluster. The expected result is again a strict JSON to be processed by \texttt{parse\_ai\_json}, then stored as \texttt{cluster.ai\_analysis} (JSONB):

\begin{lstlisting}[style=pythonstyle]
PROMPT_AI_ANALYSIS = """You are a senior clinical data analyst reviewing a
specific patient group identified through K-Prototypes population segmentation.

IMPORTANT RULES:
- Base every claim strictly on the numbers provided.
- Gender: "balanced" if 40-60%, "more female/male" if 60-70%, "predominantly
  female/male" only if >70%.
- Comparisons: use the "vs population avg" deltas to say above/below average.
- Translate vitals into plain risk language (BP >=140 = hypertensive,
  glucose >=126 = diabetic range, BMI >=30 = obese).

POPULATION CONTEXT:
- {n_patients} patients across {n_clusters} clusters
  | Silhouette: {silhouette} | Davies-Bouldin: {davies_bouldin}
- Population averages: Age {pop_avg_age} | BMI {pop_avg_bmi}
  | Systolic BP {pop_avg_systolic_bp} | Glucose {pop_avg_glucose}
  | QALY {pop_avg_qaly} | DALY {pop_avg_daly}
  | Active conditions {pop_avg_conditions} | Meds {pop_avg_meds}
  | Procedures {pop_avg_procedures}
  | Days since last encounter {pop_avg_days_since}
- Population disease rates: Diabetes {pop_diabetes_pct}%
  | Hypertension {pop_hypertension_pct}% | Heart {pop_heart_pct}%
  | Cancer {pop_cancer_pct}% | Polypharmacy {pop_polypharmacy_pct}%

THIS CLUSTER - "{title}" ({cluster_type}, {count} patients, {pct}% of population):
DEMOGRAPHICS:
- Mean age: {avg_age} (vs pop {pop_avg_age})
- Gender: F {female_pct}% / M {male_pct}%  | Marital mode: {dominant_marital}
VITALS & QUALITY OF LIFE:
- BMI {avg_bmi} (vs {pop_avg_bmi}) | Systolic BP {avg_systolic_bp} (vs {pop_avg_systolic_bp})
- Glucose {avg_glucose} (vs {pop_avg_glucose})
- QALY {avg_qaly} (vs {pop_avg_qaly}) | DALY {avg_daly} (vs {pop_avg_daly})
UTILIZATION & BURDEN:
- Active conditions {avg_conditions} (vs {pop_avg_conditions})
- Medications {avg_meds} (vs {pop_avg_meds})
- Procedures {avg_procedures} (vs {pop_avg_procedures})
- Days since last encounter {avg_days_since} (vs {pop_avg_days_since})
- Dominant encounter class: {dominant_encounter}
- Dominant careplan: {dominant_careplan}
DISEASE PREVALENCE:
- Diabetes {diabetes_pct}% (vs pop {pop_diabetes_pct}%)
- Hypertension {hypertension_pct}% (vs pop {pop_hypertension_pct}%)
- Heart disease {heart_disease_pct}% (vs pop {pop_heart_pct}%)
- Cancer {cancer_pct}% (vs pop {pop_cancer_pct}%)
- Polypharmacy {polypharmacy_pct}% (vs pop {pop_polypharmacy_pct}%)
- Smokers: current {current_smoker_pct}% / former {former_smoker_pct}%
  / never {never_smoker_pct}%

TASK: Analyse this patient group for TWO audiences:
1. CLINICIANS - 2-3 sentences anchored on dominant careplan + top prevalence
   deltas, the main risks, and the most actionable care gaps.
2. GENERAL PUBLIC - 1-2 sentences in plain everyday language, no medical terms.

Respond ONLY with valid JSON:
{{
  "risk_level": "high or moderate or low",
  "medical_interpretation": "...",
  "layperson_summary": "...",
  "key_characteristics":   ["...", "...", "...", "..."],
  "priority_interventions": ["...", "...", "..."],
  "care_recommendations":  ["...", "...", "..."]
}}
"""
\end{lstlisting}

\noindent The combination of both prompts makes up the AI part of the clusters dashboard; the first prompt is executed once per training to generate the name + color + one-liner that will be displayed on each individual card, whereas the second prompt is generated lazily after each user click to generate the actual medical interpretation, the layman’s version, and the action items list.

\subsubsection*{Libraries}

The principal libraries used are \texttt{kmodes} (algorithm), \texttt{scikit-learn} (\texttt{StandardScaler}, \texttt{MinMaxScaler}, \texttt{RobustScaler}, \texttt{PCA}, evaluation metrics), \texttt{prince} (\texttt{FAMD}), \texttt{matplotlib} and \texttt{seaborn} (plots), \texttt{pandas}/\texttt{numpy} (feature pipeline), \texttt{SQLAlchemy} (read/write paths), and the OpenAI Python SDK behind \texttt{ask\_ai} for the cluster-narrative step.

\subsection{PrefixSpan Pipeline Implementation}

The pipeline for pattern mining can be found in \texttt{app/\allowbreak clustering\_\allowbreak algorithms/\allowbreak prefix\_\allowbreak span/\allowbreak prefix\_\allowbreak span.py}, where it can run as a stand-alone script, either cluster-wide or per cluster depending on the command line option \texttt{--n-clusters N}. The mining itself is done using the \texttt{prefixspan} module, whereas everything else is about data preparation, scoring, deduplication, patient assignment, and database storage. The present script is a part of the ml-batch service. All numerical parameters are grouped together at the top of the file to make their trade-offs clear. All of them are presented and explained below:

\begin{lstlisting}[style=pythonstyle]
MIN_SUPPORT_PCT       = 0.01      # absolute floor: max(2, 1% of patients)
MIN_PATTERN_LENGTH    = 3
MAX_PATTERN_LENGTH    = 6
MAX_EVENTS_PER_PATIENT = 25
TOP_K                 = 500
WINDOW_DAYS           = 730       # 2 years of recent events per patient
OVERMINE_MULTIPLIER   = 20        # mine 10000, keep 500 after rerank
JACCARD_THRESHOLD     = 0.80
MAX_PER_PREFIX        = 6
MAX_PER_SUFFIX        = 20
TYPE_DIVERSITY_BONUS  = 0.3
MAX_DENTAL_PATTERNS   = 15

\end{lstlisting}

\subsubsection*{Event Extraction}

A single \texttt{UNION ALL} over \texttt{condition}, \texttt{medication}, \texttt{procedure}, and \texttt{careplan} results in a combined event stream. The rows' \texttt{description} is prefixed with a tag denoting the source table (\texttt{C:}, \texttt{M:}, \texttt{P:}, \texttt{CP:}) so that when the same string appears in two tables, it is still counted as two unique strings:

\begin{lstlisting}[style=sqlstyle]
SELECT patient_id, start_date AS event_date, 'C:' || TRIM(description) AS item
FROM condition
WHERE description IS NOT NULL AND start_date IS NOT NULL
UNION ALL
SELECT patient_id, start_date, 'M:' || TRIM(description) FROM medication
UNION ALL
SELECT patient_id, start_date, 'P:' || TRIM(description) FROM "procedure"
UNION ALL
SELECT patient_id, start_date, 'CP:' || TRIM(description) FROM careplan;

\end{lstlisting}

\subsubsection*{Sequence Construction and Three Carefully Justified Filters}

The deterministic sort order used by the \texttt{build\_sequences} function is \texttt{(patient\_id, event\_date, \_pref\_rank, item)} with \texttt{kind="mergesort"} for stability, meaning ties within the same day will always be broken by the following precedence order: \texttt{C} \(<\) \texttt{M} \(<\) \texttt{P} \(<\) \texttt{CP}. The three filter functions execute in a pipeline and each is needed to avoid a particular pitfall of the initial prototypes.

\paragraph{(1) Substring blacklist.} Several substrings that are common forms of garbage data are discarded:

\begin{lstlisting}[style=pythonstyle]
ITEM_BLACKLIST_SUBSTRINGS: frozenset[str] = frozenset({
    "assessment of health and social care needs",
    "medication review due",
    "medication reconciliation",

})
\end{lstlisting}

\noindent The strings mentioned here are present in almost every episode in the Synthea database and are not patterns but simply workflow indicators. All top-K results had one of these strings as their leading item in the absence of the blacklist. This meant that the actual important patterns were being overshadowed. The frozenset matches using substring matching (\texttt{str.contains(sub, regex=False)}), as Synthea often adds a code to the string.

\paragraph{(2) Recency window.} All patient interactions that occurred more than \texttt{WINDOW\_DAYS=730} days before that patient's most recent interaction are dropped from the sample using \texttt{groupby(patient\_id).transform("max")} so that the window is calculated independently for each patient without leaking information at the dataset level:

\begin{lstlisting}[style=pythonstyle]
maxd  = events.groupby("patient_id", sort=False)["event_date"].transform("max")
cutoff = maxd - pd.Timedelta(days=WINDOW_DAYS)
events = events[events["event_date"] >= cutoff]
\end{lstlisting}

\noindent The maximum value of 730 days (2 years) ensures clinical relevance, since the sequence of treatments that occurred many years prior would not have any clinical significance and would only make the sequences longer than needed. Additionally, the cutoff is anchored to the last event of each individual, no matter how many years ago that person died. Therefore, the last 2 years of each patient’s medical history are included.

\paragraph{(3) Consecutive duplication removal and maximum timeline length.} In case there are two consecutive identical elements for one patient, the second one is dropped based on comparisons of \texttt{patient\_id} and \texttt{item} with \texttt{shift()}. \texttt{MAX\_EVENTS\_PER\_PATIENT=25} caps each timeline via \texttt{groupby().tail(25)}.

\begin{lstlisting}[style=pythonstyle]
same_pid  = events["patient_id"].eq(events["patient_id"].shift())
same_item = events["item"].eq(events["item"].shift())
events    = events[~(same_pid & same_item)]
events    = events.groupby("patient_id", sort=False, group_keys=False).tail(MAX_EVENTS_PER_PATIENT)
\end{lstlisting}

\noindent Without duplicate elimination, there will be a huge number of duplicates for the \texttt{M:Lisinopril} row that a patient will get when using some medicine regularly; therefore, PrefixSpan can connect those rows and form the trivial pattern \texttt{M:Lisinopril \(\rightarrow\) M:Lisinopril \(\rightarrow\) M:Lisinopril}. The scanning time required without the upper bound on the timeline length will be \(O(\text{tail}^2)\), where each super-utilizer projection table may have hundreds of rows.

\subsubsection*{Encoding and Mining Loop}

The sequences are transformed into integer IDs with the help of the one-pass vocabulary builder (\texttt{\_encode}). This way, we ensure that the inner workings of the \texttt{prefixspan.PrefixSpan} function rely on numerical data and not on textual strings. This is definitely the single most important speedup in the entire process pipeline. The package is invoked with the closed top-\(k\) operator, configured to mine a candidate pool \(20\) times larger than the desired output size:

\begin{lstlisting}[style=pythonstyle]
ps         = PrefixSpan(int_seqs)
ps.minlen  = MIN_PATTERN_LENGTH        # 3
ps.maxlen  = MAX_PATTERN_LENGTH        # 6
pool_size  = TOP_K * OVERMINE_MULTIPLIER  # 10000
raw        = ps.topk(pool_size, closed=True)
abs_minsup = max(2, int(n_patients * MIN_SUPPORT_PCT))   # 1%
raw        = [(sup, pat) for sup, pat in raw if sup >= abs_minsup]

\end{lstlisting}

\noindent Over-Mining: This is the optimization technique applied to allow all downstream filters to be aggressive in pruning, without running into the danger of mining too few patterns for consideration. For instance, when we call \texttt{top\_k=500}, and rescoring, deduplication, prefix/suffix capping, and dental capping filter out roughly half of the number, leaving us with 250 candidates which is well below what the dashboard requires. Mining 20 times more and applying filtration afterwards enables the filters to be aggressive in pruning the near-duplicates, without necessarily yielding a thin set of patterns. Using \texttt{closed=True} allows the output to include only closed patterns.

The minimum support threshold is set to \texttt{int(n\_patients * MIN\_SUPPORT\_PCT)} (\(1\%\)) with a hard lower bound of 2. Applying the minimum \(1\%\) makes it possible to bring out even very small and niche-specific pathways (such as a maternal pathway in an overwhelmingly non-pregnant cohort), while at the same time preventing a single noisy patient from elevating a pattern.

\subsubsection*{Geometric-Mean-Lift Scoring}

The new evaluation for each candidate will be based on the geometric mean lift (GML) measure that calculates the ratio between the actual support and the product of the marginal frequencies of the events in the pattern:

\begin{lstlisting}[style=pythonstyle]
log_item_freqs = np.log(np.maximum(single_support[unique_items], 1)) - log_n
raw_gml        = np.exp(np.log(sup) - log_n - np.mean(log_item_freqs))
gml            = raw_gml * (1.0 + TYPE_DIVERSITY_BONUS * (n_types - 1))

\end{lstlisting}

Furthermore, there is a multiplicative factor (\texttt{TYPE\_DIVERSITY\_BONUS}), determined by the number of distinct event types (C/M/P/CP) appearing in the pattern; therefore, a pattern in the format \texttt{C: \(\rightarrow\) M: \(\rightarrow\) P:} will have a higher ranking than an ordinary sequence consisting of events from one single domain despite having a comparable support. Patterns with three consecutive events of the same type are filtered from the candidate list beforehand.

\subsubsection*{Closedness and Jaccard Deduplication}

Finally, a strictly closed set of patterns will be generated by removing a pattern \(p\) from the candidate list when a pattern \(p'\) longer than \(p\) with the same support contains all the elements of \(p\) within its item set, where \(p\) appears as a subsequence of \(p'\). This will be performed in buckets with the keys corresponding to the support values to avoid quadratic scanning over the whole candidate list; thus, achieving the time complexity of \(O(\text{bucket}^2)\). Deduplication will occur at the following stage:

\begin{lstlisting}[style=pythonstyle]
def _jaccard(a, a_len, b, b_len):
    inter = len(a & b)
    denom = a_len + b_len - inter
    return inter / denom if denom else 0.0

if all(_jaccard(pat_set, pat_len, s, s_len) < JACCARD_THRESHOLD
       for s, s_len in chosen_sets):
    chosen.append(entry)

\end{lstlisting}

Two caps at the end ensure that the impact of any single starting or ending element does not unfairly affect the outcome (\texttt{MAX\_PER\_PREFIX}=6, \texttt{MAX\_PER\_SUFFIX}=20). Additionally, there is another cap to limit the number of majority dental pathways among the top-\(k\) pathways (\texttt{MAX\_DENTAL\_PATTERNS}=15). This cap is critical due to the fact that Synthea has a standard routine that involves applying a fluoride product and other dental treatments (\texttt{P:Application of fluoride product} \(\rightarrow\) \texttt{P:Dental cleaning} \(\rightarrow\) \texttt{P:Dental examination} etc.), which leads to many near-duplicate patterns. Without such a cap, the 50 most common pathways for the cohort would be completely filled with dental routines, whereas cardiometabolic and respiratory patterns would not appear on the dashboard. Detection of majority dental pathways is based on keywords:

\begin{lstlisting}[style=pythonstyle]
DENTAL_KEYWORDS: frozenset[str] = frozenset({
    "dental", "fluoride", "gingiv", "supragingival", "subgingival",
    "oral health", "tooth", "teeth",
})

\end{lstlisting}

\subsubsection*{Libraries}

The pattern discovery process relies on the constrained closed top-\(k\) approach of the \texttt{prefixspan} algorithm. The data manipulation tasks like filtering and preprocessing are handled via \texttt{pandas} library, whereas \texttt{numpy} is utilized for math-based functions such as geometric mean lift scores, diversity scores, and support count calculation. Database access and persistence-related functionalities are managed via \texttt{SQLAlchemy} (initialized via \texttt{python-dotenv}), where \texttt{ai\_analysis} is intentionally set as \texttt{NULL} for lazy loading of LLM analysis. Finally, the basic python libraries serve the purpose of providing infrastructure functionality such as \texttt{collections.defaultdict} indexing patient-to-pattern mapping, \texttt{argparse} parsing of scoping flags, and \texttt{pathlib}.

\subsection{LLM SQL Agent Implementation}

The SQL agent itself can be found in \texttt{app/service/lllm\_sql\_agent\_service/}, inside the \texttt{analytics microservice}, and is the gateway through which untyped natural-language questions can be posed to the database. The SQL agent is built on \textbf{LangGraph}, which represents the workflow as a directed graph of node functions operating on a typed shared state, compiled once at import time and run per request by invoking \texttt{graph.invoke(initial\_state, \{"recursion\_limit": 20\})}. A schematic of the graph can be seen in Figure~\ref{fig:ch5_sql_agent_graph}. Each node accepts the current state as input, issues at most one LLM call (using \texttt{llm\_models.ask\_ai}, a wrapper around the OpenAI Python SDK, configured with \texttt{timeout=60s} and \texttt{max\_retries=2}), and emits a partial state update. The entire database context injected into every prompt - 15 tables, all column types, each categorical value set, every foreign-key mapping, and the \texttt{cluster\_patient.features} JSONB access convention - is materialised as the string constant \texttt{DB\_SCHEMA} in \texttt{db/db\_data.py}.

\begin{figure}[H]
\centering
\includegraphics[width=0.90\linewidth]{figs/chapter_5/sql_graph.png}
\caption{LangGraph state machine for the SQL agent}
\label{fig:ch5_sql_agent_graph}
\end{figure}

First, \texttt{process\_input} classifies the user message into one of three categories: \texttt{DB\_\allowbreak RETRIEVAL} (rewritten against the previous 5 chat turns), \texttt{MEDICAL\_\allowbreak GENERAL} (health-related query that does not require the database), or \texttt{NOT\_\allowbreak MEDICAL} (any other query, treated as out of scope). For database-bound queries, \texttt{generate\_query} drafts a SELECT-only SQL statement, \texttt{validate\_query} applies a static safety filter followed by a second-LLM semantic audit, and finally \texttt{execute\_query} runs the validated SQL through the request-scoped SQLAlchemy session, truncating the raw result to 3000 characters before \texttt{PROMPT\_RESULT} summarises it in the user's original language. On failure at either step, the error message is fed back into the next round of query generation. The cycle is bounded by \texttt{MAX\_RETRIES=3}. If the calling user has role \texttt{DOCTOR}, then \texttt{DOCTOR\_\allowbreak SCOPE\_\allowbreak CLAUSE}, which carries the \texttt{provider\_id} of that particular doctor, is injected into both the generation and validation prompts, and a post-execution textual check refuses to return the answer if the final SQL does not reference that provider id.

A single end-to-end prompt would be inadequate for three reasons addressed by the architecture. First, a one-shot prompt could never guarantee safety: the static regular-expression filter (rejecting \texttt{DROP/DELETE/TRUNCATE/ALTER/INSERT/UPDATE/CREATE/GRANT/REVOKE/EXEC} and requiring the statement to start with \texttt{SELECT} or \texttt{WITH}) is deterministic Python that no amount of prompt injection could convince to do otherwise. Second, even when the SQL is safe, the model can hallucinate columns, omit the \texttt{NOT is\_outlier} filter required for cluster queries, or use \texttt{COUNT(*)} when it should have \texttt{COUNT(DISTINCT patient\_id)}; the validator acts as a second pair of eyes independent of the context biases of the generator. Third, on failed validation or failed execution, the agent does not merely retry blindly; it regenerates with the exact previous SQL and the exact cause of the failure injected as a \texttt{retry\_context} block in the prompt. The synergy of the following three properties is what gives the agent the robustness to self-correct: (a) schema injection as static context, (b) separation of generation and auditing into cooperating LLM calls, and (c) feedback-driven regeneration on the retry path.

\subsubsection*{Database Context Asset}

Every prompt requiring reasoning about the database is grounded against one static resource, \texttt{DB\_\allowbreak SCHEMA}, declared in \texttt{db/db\_data.py}. It is a meticulously curated description of all 15 tables (4 infrastructure, 7 clinical, 2 for sequence mining, and 3 for clustering) that enumerates every column, type, foreign-key relationship, and JSONB access convention. \texttt{DB\_\allowbreak SCHEMA} is injected verbatim into \texttt{PROMPT\_\allowbreak TRANSLATE}, \texttt{PROMPT\_\allowbreak GENERATE\_\allowbreak SQL}, and \texttt{PROMPT\_\allowbreak VALIDATE\_\allowbreak SQL}:

\begin{lstlisting}[style=pythonstyle]
DB_SCHEMA = """
=== HOSPITAL ANALYTICS DATABASE SCHEMA ===

OVERVIEW:
- 15 tables total: 4 infrastructure (patient, organization, provider, encounter)
  + 7 clinical (observation, condition, medication, procedure, immunization,
    allergy, careplan) + 2 sequence-mining + 3 clustering ...
- The central entity is Patient. Encounters link Patients to Providers ...

--- TABLES ---
TABLE: patient
  - patient_id (str, PK) - UUID. Never treat as integer.
  - gender (str) - Values: 'M' (Male), 'F' (Female). Do NOT use 'Male'/'Female'.
  - ...
TABLE: encounter
  - encounter_class - 'inpatient', 'ambulatory', 'outpatient', 'wellness', ...
  - ...
"""
\end{lstlisting}

\subsubsection*{Prompts Used by the SQL Agent}

\paragraph{\texttt{PROMPT\_TRANSLATE} - input classifier and rewriter.} Classifies all messages as belonging to one of three categories, transforming database queries to complete English sentences utilizing past conversations and the schema to resolve term ambiguity.

\begin{lstlisting}[style=pythonstyle]
PROMPT_TRANSLATE = """
You classify and rewrite user messages for a hospital analytics assistant.
User message: "{initial_prompt}"
{chat_context}
Schema reference (for term disambiguation only, do NOT copy into output):
{db_schema}

Classification rules - return EXACTLY one of:
1. DB_RETRIEVAL  -> a single standalone English sentence describing what data to retrieve.
2. MEDICAL_GENERAL: <question>
3. NOT_MEDICAL
Output ONLY the chosen output. No markdown, no quotes, no extra explanation.
"""
\end{lstlisting}

\paragraph{\texttt{PROMPT\_\allowbreak MEDICAL\_\allowbreak ANSWER} - non-DB medical answer.} Triggered whenever the classifier produces \texttt{MEDICAL\_\allowbreak GENERAL}. Ensures the answer stays in the user's language and appends a safety nudge for personal health-related questions.

\begin{lstlisting}[style=pythonstyle]
PROMPT_MEDICAL_ANSWER = """
You are a medical assistant for a hospital analytics platform.
{chat_context}
Question: {question}
Rules:
- Reply in the SAME language the user wrote in.
- Be accurate and concise: 2-5 sentences or a short bulleted list.
- If the question implies a personal health situation, advise to consult a doctor.
- Never mention SQL, databases, or technical system internals.
"""
\end{lstlisting}

\paragraph{\texttt{PROMPT\_\allowbreak GENERATE\_\allowbreak SQL} - SELECT generator.} Writes the SQL query based on the schema, optionally including the scope of a role, and any feedback from past failures. Incorporates every idiosyncrasy of the project: \texttt{LIMIT}, \texttt{ILIKE} for free text, exact enums for categorical codes, \texttt{COUNT(DISTINCT patient\_id)}, JSONB accessors for \texttt{cluster\_patient.features}, the mandatory \texttt{NOT is\_outlier} filter, and the cluster/sequence join rules.

\begin{lstlisting}[style=pythonstyle]
PROMPT_GENERATE_SQL = """
You are a PostgreSQL expert writing a SELECT query against a hospital analytics database.
Question: "{prompt}"
Schema: {db_schema}
{scope_clause}{retry_context}

Output rules:
1. Return ONLY the raw SQL. No markdown, no prose.
2. ONLY SELECT, or WITH ... SELECT. Never INSERT/UPDATE/DELETE/DROP/...
3. Select only the columns needed; never SELECT *.
4. Add LIMIT {max_rows} unless the user asks for "all" or it's aggregate-only.
...
Filtering rules:
7. Free-text fields - use ILIKE with %...% wildcards.
8. Categorical codes - use `=` with exact values (gender 'M'/'F', encounter_class ...).
...
Clustering rules:
13. cluster_patient: always filter `NOT is_outlier` unless the user asks about outliers.
14. cluster_patient.features is JSONB. Use features->>'key' / (features->>'key')::numeric.
...
"""
\end{lstlisting}

\paragraph{\texttt{PROMPT\_\allowbreak VALIDATE\_\allowbreak SQL} - second-LLM auditor.} Independent semantic verification. Responds \texttt{YES} on line 1 or \texttt{NO} along with a one-sentence rationale on line 2, which serves as feedback for the subsequent attempt.

\begin{lstlisting}[style=pythonstyle]
PROMPT_VALIDATE_SQL = """
You are a SQL QA auditor. Decide if the SQL is safe AND correctly answers the request.
User Request: "{prompt}"
Generated SQL: {sql_query}
Schema: {db_schema}
{scope_clause}

Validation checklist (reject if ANY fails):
1. Must be a single SELECT or WITH ... SELECT.
2. Every referenced table and column must exist in the schema.
3. JOIN conditions must use documented foreign-key relationships.
4. Query logic must answer THIS specific question.
5. Cluster questions: reject if `NOT is_outlier` is missing.
6. Patient counts across clinical tables: reject COUNT(*) where COUNT(DISTINCT patient_id) is needed.
7. If a MANDATORY ROLE SCOPE is present: reject if the provider_id filter is missing.

Do NOT flag: ILIKE vs =, aliasing, ordering, whitespace, LIMIT on aggregate-only.

Output format (strict):
Line 1: YES or NO
Line 2 (only if NO): one-sentence explanation.
"""
\end{lstlisting}

\paragraph{\texttt{PROMPT\_\allowbreak RESULT} - natural-language summariser.} Runs once after a successful execution. Phrases the answer as a data analyst would, never as a DBA.

\begin{lstlisting}[style=pythonstyle]
PROMPT_RESULT = """
Summarise database results for the user in natural language.
User question: "{initial_prompt}"
Database result (list of row dicts): {result}

Rules:
- Reply in the SAME language the user wrote in.
- Empty result -> say plainly that no matching records were found.
- Single value -> one sentence with the number and unit.
- Multi-row -> short bulleted/numbered list; truncate to top ~10.
- Never mention SQL, queries, JOINs, tables, columns, or technical internals.
- Do not invent or extrapolate values beyond what the result contains.
"""
\end{lstlisting}

\paragraph{\texttt{DOCTOR\_\allowbreak SCOPE\_\allowbreak CLAUSE} - role-based scope injection.} Conditionally prepended to both \texttt{PROMPT\_\allowbreak GENERATE\_\allowbreak SQL} and \texttt{PROMPT\_\allowbreak VALIDATE\_\allowbreak SQL} depending on whether the call was initiated by a doctor or not. A post-execution textual check then refuses the answer if the final SQL does not reference the doctor's \texttt{provider\_id}.

\begin{lstlisting}[style=pythonstyle]
DOCTOR_SCOPE_CLAUSE = """
MANDATORY ROLE SCOPE - DOCTOR (provider_id = '{provider_id}'):
- Every SQL query MUST restrict results to encounters where provider_id = '{provider_id}'.
- For patient questions, restrict to: patient.patient_id IN
  (SELECT DISTINCT patient_id FROM encounter WHERE provider_id = '{provider_id}').
- Never expose other doctors' patients, encounters, or aggregate counts.
"""
\end{lstlisting}